\clearpage
\section*{Anexo 4-P --- Tecnologías, soporte y actualización}
\phantomsection\label{fuente-19-anexo-4-1-p-tecnologias-soporte}
\phantomsection
\label{anx:P}
\addcontentsline{toc}{section}{Anexo 4-P --- Tecnologías, soporte y actualización}

Este anexo relaciona los nombres funcionales con los servicios técnicos y declara, para cada tecnología, su versión, su soporte y su política de actualización.
\subsection*{Detalle de selección por capa}

La selección siguiente organiza las tecnologías por responsabilidad arquitectónica y explica su función, compatibilidad y continuidad durante el contrato.

\begin{itemize}
\item \textbf{Backend (12 módulos).} Laravel 13 sobre PHP 8.5, con Composer y el archivo de versiones de dependencias, organizado en contextos M1–M12. PostgreSQL/PostGIS conserva los datos geográficos. M4 usa repositorios espaciales con SQL parametrizado vía PDO/Query Builder y pruebas de consultas. Laravel 13 requiere PHP 8.3 o superior y admite PHP 8.5 (Laravel, 2026). Ambos componentes se actualizarán a versiones con soporte durante los 56 meses.

\item \textbf{Frontend web.} Angular 22 con Tailwind CSS y TypeScript 6.0 compatible con su matriz oficial. La integración con Keycloak se implementa mediante OIDC y se mantiene junto con las dependencias del cliente. Las versiones deben actualizarse durante el contrato conforme a sus ciclos de soporte.

\item \textbf{Apps de campo (preventa y reparto).} Kotlin Android nativo. Nativo del parque Zebra/Android, escáner GS1 vía Zebra DataWedge, SQLite/Room offline, acceso nativo a GPS, POS y térmica.

\item \textbf{Borde y continuidad del acceso (Capa 2).} CloudFront, AWS WAF, AWS Shield Advanced y ALB en nube, con control de acceso local. En los CD, Starlink fijo permanece encendido con IPsec establecido y BGP de menor preferencia tras fibra y LTE. En los cross-docking es el camino principal, con LTE de dos proveedores como respaldo y conmutación automática en menos de 30 segundos (ADR-02), sin sustituir la autonomía del servicio local.

\item \textbf{Puerta de enlace de servicios (Capa 3).} Amazon API Gateway para publicar las APIs y aplicar controles de entrada. La configuración debe cubrir la validación de identidad de Keycloak, las cuotas y los límites de solicitudes, conforme a la modalidad de API seleccionada y al flujo descrito en la capa de puerta de enlace de este apartado.

\item \textbf{Base de datos transaccional on-premise.} PostgreSQL + PostGIS por sitio con el mismo perfil local de bodega. Cada sitio es autoridad de sus movimientos y N-04/N-05 consolida stock y reservas en nube.

\item \textbf{Base de datos nube (OLTP y recuperación).} Amazon Aurora PostgreSQL, con despliegue Multi-AZ y una ventana propuesta de recuperación a un punto en el tiempo (PITR) de 35 días. Los tiempos de conmutación y restauración se verifican mediante pruebas frente a los objetivos RTO/RPO establecidos.

\item \textbf{Ingesta IoT / frío.} Amazon DynamoDB e IoT Core reciben las lecturas. Greengrass corre en tres gateways Moxa UC-8200 o equivalentes: dos en Talca, que leen ambas cámaras por Modbus TCP para evitar un punto único de falla, y uno en Concepción. Leen sensores por Modbus, bloquean el despacho ante una excursión crítica y sostenida en menos de 5 segundos y conservan 24 horas sin enlace. Los termógrafos de los camiones registran toda la ruta en memoria interna y difunden por BLE la temperatura. El terminal alerta la excursión en la cabina, reenvía la alerta por la red celular, de inmediato con cobertura o al recuperarla, y descarga el registro completo al volver. Escrituras serverless con TTL nativo (raw 30 días).

\item \textbf{Serie consolidada (OLAP).} S3 Parquet + AWS Glue + Redshift Serverless. OLAP por diseño: DynamoDB raw → Glue → S3 Parquet → Redshift, sin motor de series adicional.

\item \textbf{Caché y datos de turno.} Redis acelera las lecturas centrales. No se exige Redis en cada sitio. Las APIs entregan una copia cifrada para SQLite/Room y conservan por separado la cola de escrituras sin confirmar. La identidad usa credenciales de 8 horas en bodega y 14 horas en terreno. Su caché local de solo lectura tiene TTL de 24 horas en VM-05, VM-C03 y cada E-01, junto al verificador local de relevos (ADR-06).

\item \textbf{Mensajería asíncrona (Capa 5).} RabbitMQ en on-premise mediante adaptador AMQP \texttt{php-amqplib}, SQS FIFO con sobre JSON y consumidor PHP para reconciliación, colas SQS separadas para trabajos Laravel, y SNS para difusión y alertas en nube. Un único planificador Laravel programa las tareas periódicas. \textbf{Dos procesos del servicio de integración con el ERP en VM-04 consumen RabbitMQ local y, por salida, la cola FIFO de solicitudes al ERP, ERP integrado exclusivamente vía ACL, Hub EDI GS1 centralizado en nube (EANCOM, GS1 XML, EPCIS), AS2 como transporte.}

\item \textbf{Analítica / BI.} S3 Data Lake, Redshift Serverless y Glue ETL. Consultas históricas sin degradar el procesamiento transaccional.

\item \textbf{Objetos / documentos.} Amazon S3 con Intelligent-Tiering y Object Lock.

\item \textbf{IAM / Identidad (Capa 7).} Keycloak (OIDC, SAML, SSO, MFA) como IdP maestro en nube, con caché local de solo lectura de TTL 24 horas y verificador local con manifiesto firmado y PIN personal en VM-05, VM-C03 y cada E-01.

\item \textbf{Gestión de secretos (Capa 7).} AWS Secrets Manager y SSM Parameter Store con rotación automática. Cuenta de emergencia fuera de banda, con doble autorización, registro de uso y rotación posterior de credenciales.

\item \textbf{Observabilidad (Capa 8).} OpenTelemetry para PHP/Laravel y colectores ADOT on-premise con almacenamiento temporal en disco de 24 horas, plataforma única Amazon CloudWatch para logs (12 meses en línea + 24 en archivo), métricas (13 meses), trazas y tableros (ADR-14).

\item \textbf{Gestión de dispositivos (RT-03.18, Bases Técnicas Transversales, cap. 3, p. 9).} Gestión centralizada mediante MDM (Android Enterprise / Zebra DNA, SaaS), con inventario, configuración, actualización de firmware y aplicaciones, bloqueo y borrado remoto del parque móvil.

\item \textbf{Contenedores / orquestación.} Imagen PHP 8.5 con servidor HTTP/PHP-FPM para APIs y procesos PHP CLI separados para colas y programación. ECS Fargate aloja los perfiles centrales. Talca y Concepción ejecutan el perfil local de bodega sobre máquinas virtuales Proxmox, y Docker Compose ejecuta el mismo perfil local de bodega en los equipos E-01 de los cross-docking. El código y las migraciones de esquema proceden del mismo artefacto versionado.

\item \textbf{Infraestructura como Código.} Terraform administra recursos mediante proveedores y estados separados por ambiente. Ansible configura hosts. CDK no administra los mismos recursos: cada activo tiene un único propietario de infraestructura como código.

\item \textbf{CI/CD.} GitLab CI como orquestador + AWS CodeBuild para construcción hermética con procedencia SLSA 3, la auditoría de dependencias, PHPUnit, PHPStan/Larastan y Laravel Pint verifican dependencias, comportamiento, análisis estático y formato. \texttt{swagger-php} genera OpenAPI 3.1 desde atributos PHP. Los esquemas AsyncAPI 2.6 se validan y publican en la misma puerta de calidad.
\end{itemize}

\paragraph{Administración del parque por familia de dispositivo}

Mantiene inventario central, configura, actualiza firmware y aplicación, bloquea y borra de forma remota el parque de terreno y bodega (RT-03.18, Bases Técnicas Transversales, cap. 3, p. 9). Las reservas se entregan precargadas (Bases Técnicas Transversales, cap. 8, p. 19). El resto de los dispositivos se gestiona desde la consola de su fabricante e ingresa al mismo inventario: las impresoras ZQ620 Plus y ZT411 con Zebra Printer Profile Manager Enterprise sobre Link-OS. Los terminales PAX A920Pro Duo Premium con PAXSTORE, que permite configurar, actualizar, bloquear y retirar cada equipo. Los gateways IoT con la gestión de flota de AWS IoT Greengrass, y los termógrafos con ELA Device Manager por NFC, mientras la aplicación de reparto supervisa su batería y su registro en cada ruta. Las balanzas y los módulos de sensores no tienen sistema operativo administrable: se registran en el inventario y su firmware se actualiza en la mantención (RT-08.14, Bases Técnicas Transversales, cap. 8, p. 19).

\paragraph{Correspondencia entre nombres funcionales y técnicos}
El inventario de servicios AWS comprende CloudFront, AWS WAF, AWS Shield Advanced, ALB, Network Load Balancer, API Gateway, Verified Access, Aurora, ElastiCache, DynamoDB, S3, Redshift Serverless, ECS Fargate, Lambda (solo como autorizador), Route 53, Secrets Manager, Systems Manager, KMS, IAM, Organizations, IAM Identity Center, CloudWatch, Transit Gateway, NAT Gateway, VPC Endpoints, Site-to-Site VPN, SQS, SNS, AWS Backup, GuardDuty, Security Hub, Security Lake, Macie, Inspector, Config, Database Migration Service, Glue, QuickSight, Elastic Container Registry, CodeBuild, Control Tower, CloudTrail e IoT Core+Greengrass.

El perfil local de bodega se identifica técnicamente como \texttt{wms\_only}; el servicio de integración con el ERP, como \texttt{erp-sync}; el agente de transferencia del sitio, como \textit{shipper}; y la bandeja transaccional de salida, como \textit{outbox}. La separación interna de dominio, aplicación, infraestructura y acceso HTTP corresponde a \texttt{Domain}, \texttt{Application}, \texttt{Infrastructure} y \texttt{Http}, respectivamente. Composer fija las versiones de dependencias en \texttt{composer.lock} y las audita mediante \texttt{composer audit}. El versionado mayor, menor y de corrección corresponde a \texttt{major.minor.patch}. En 4.1, el registro versionado de parámetros corresponde a \texttt{mae\_parametro} y \texttt{mae\_parametro\_version}; el registro de asignaciones, a \texttt{gob\_asignacion}; el registro de auditoría, a \texttt{gob\_auditoria}; y la plantilla de notificación, a \texttt{not\_plantilla} con su \texttt{cuerpo\_ref}. Estas equivalencias conservan los nombres de implementación utilizados por los contratos, los ADR y el despliegue de 4.2.

El inventario distingue versión de referencia de una imagen exacta de producción. Cada liberación fija parches, hashes y compatibilidad en los archivos de bloqueo y SBOM. Las versiones futuras se aprueban mediante pruebas antes de promoverse. La hoja de ruta cubre los 56 meses mediante actualizaciones, no suponiendo soporte de una misma versión durante todo el contrato (Bases Técnicas Transversales, numeral 1.6, pp. 4--5).

\begin{tablalafrox}{Ciclo de soporte del núcleo lógico}{tab:soporte-logico}{F{0.16} F{0.15} F{0.23} Y}{\cab{Producto} & \cab{Referencia} & \cab{Fin de soporte publicado} & \cab{Plan de actualización}}
Laravel & 13.x & Seguridad: 17-03-2028. & Evaluación anual, migrar al menos 6 meses antes del fin. \\
PHP & 8.5.x & Seguridad: 31-12-2029. & Actualización compatible antes del fin, ensayo de APIs y colas. \\
PostgreSQL & 16.x & 09-11-2028. & Ensayo de migración mayor y reversión 6 meses antes. \\
Angular & 22.x & LTS: junio de 2028, fecha orientativa del fabricante. & Revisión semestral y actualización anual por compatibilidad. \\
TypeScript & 6.0.x & Sin fecha contractual independiente. & Versión compatible con Angular, archivo de bloqueo y pruebas. \\
RabbitMQ & Rama 4.x con soporte comunitario vigente al liberar. & 4.3.x, comunidad: 31-01-2027. & Se instala la rama vigente y se actualiza antes de su fin. No se presume licencia comercial. \\
\end{tablalafrox}
{\fontsize{9pt}{11pt}\selectfont\raggedright Fuente: elaboración a partir de Laravel (2026), PHP (2026), PostgreSQL (2026), Angular (2026a, 2026b) y RabbitMQ (2026). Fechas consultadas el 30-09-2026. Diseño propuesto: LafroX (2026e, diseño de este apartado, pp.~\pageref*{fuente-19-anexo-4-1-p-tecnologias-soporte}--\pageref*{fuente-19-anexo-4-1-p-tecnologias-soporte-fin}).\par}

Laravel, PHP, PostgreSQL y Angular tienen horizontes inferiores al contrato: el responsable de Desarrollo registra cada trimestre su soporte y planifica la sustitución antes de vencimiento. RabbitMQ requiere una actualización temprana antes del fin comunitario de la rama de referencia, aun durante desarrollo. No se atribuye al CLIENTE un contrato de soporte comercial no contratado.

La Tabla~\ref{tab:inventario-software} reúne el inventario complementario y la matriz de navegadores, versiones y actualización que se aplica a los portales y a la consola web.

\begin{tablalafrox}{Inventario complementario y control de vigencia}{tab:inventario-software}{F{0.23} F{0.23} Y}{\cab{Componente} & \cab{Línea base lógica} & \cab{Control durante 56 meses}}
PostGIS & 3.x compatible con PostgreSQL 16. & Compatibilidad y fin de soporte de la distribución registrados antes de liberar. \\
Kotlin/Android & Rama estable compatible con parque Zebra. & Matriz OS, SDK y periféricos por modelo, MDM impide OS fuera de soporte. \\
Room/SQLite & Versión estable de AndroidX Room y SQLite incorporada. & Parches fijados por Gradle, migración de cola sin perder registros sin confirmar. \\
Keycloak & 26.x, parche soportado al liberar. & Revisión trimestral del ciclo oficial, migrar junto al adaptador OIDC. \\
Tailwind/Node/RxJS & Versiones compatibles con Angular y su toolchain. & Archivo de bloqueo y matriz de compatibilidad, Node solo construye el cliente. \\
Google Chrome --- Windows, macOS, Linux y Android & Versión estable mayor vigente y la inmediatamente anterior, con sus parches disponibles. & Registrar versiones exactas y verificar ambas contra el Baseline de Angular y los flujos de los portales y consola. \\
Microsoft Edge --- Windows, macOS y Linux & Versión estable mayor vigente y la inmediatamente anterior, con sus parches disponibles. & Registrar versiones exactas y verificar ambas contra el Baseline de Angular y los flujos de los portales y consola. \\
Mozilla Firefox --- Windows, macOS, Linux y Android & Versión estable mayor vigente y la inmediatamente anterior, con sus parches disponibles. & Registrar versiones exactas y verificar ambas contra el Baseline de Angular y los flujos de los portales y consola. \\
Apple Safari --- macOS, iOS y iPadOS & Versión estable mayor vigente y la inmediatamente anterior, en sistemas compatibles. & Registrar versión del navegador y del sistema y verificar ambas contra el Baseline de Angular y los flujos de los portales y consola. \\
Composer\slash PHPUnit\slash PCOV\slash PHPStan\slash Pint\slash PHPMD\slash PHPCPD\slash PhpMetrics\slash deptrac & Versiones compatibles con PHP 8.5 y Laravel 13. & Composer lock, análisis de licencia, audit, cobertura, complejidad, duplicación, deuda, fronteras entre módulos y regresión, no entran herramientas de prueba al runtime. \\
JUnit 5/Kover/Espresso/detekt/ktlint & Versiones compatibles con Kotlin, Android y el parque Zebra. & Cobertura de dominio y sincronización, pruebas de interfaz y controles de estilo y complejidad antes de liberar el artefacto móvil. \\
 Jest\slash Playwright\slash ESLint\slash axe-core\slash OWASP ZAP\slash Pact\slash Spectral\slash MobSF\slash k6\slash AWS Fault Injection Service  &  Versiones compatibles con Angular, Android, las APIs y los ambientes AWS.  &  Regresión y pruebas de extremo a extremo de portales, lint, accesibilidad WCAG 2.2 A/AA, DAST, contratos, análisis de la aplicación Android, carga y estrés e inyección controlada de fallas en QA y Preproducción.  \\
Symfony ExpressionLanguage & Versión mantenida compatible con PHP 8.5 y Laravel 13, fijada en Composer lock. & Motor de reglas de 4.1.3.4: registrar versión y soporte, revisar vulnerabilidades y probar resultados de reglas y reversión antes de actualizar. \\
php-amqplib/AWS SDK/swagger-php & Bibliotecas mantenidas compatibles con PHP 8.5. & Contrato AMQP, sobre JSON y generación OpenAPI bloquean cambios incompatibles. \\
OpenAS2/JVM & Rama soportada compatible con perfiles de las cadenas. & Certificados, licencia, runtime Java y pruebas MDN registrados antes de habilitar INT-08. \\
OpenTelemetry/ADOT & Versiones estables de SDK y colector. & Ensayo de recepción OTLP, correlación, buffer y exportación tras corte. \\
Terraform/Ansible & Versiones estables y providers fijados. & Terraform gobierna recursos, Ansible configura hosts. AWS CDK no gobierna los mismos recursos. \\
Servicios AWS & API/engine gestionados declarados por servicio. & Avisos de retiro trimestrales. Prueba de adaptación. Aurora conserva versión de engine y fecha de soporte del proveedor. \\
MDM/Android Enterprise & Servicio y plan compatibles con el parque. & Contrato de soporte y actualización de dispositivos durante todo el período. \\
\end{tablalafrox}
{\fontsize{9pt}{11pt}\selectfont\raggedright Fuente: elaboración a partir de decisiones de arquitectura lógica y política de configuración de la propuesta, compatibilidad de navegadores según Angular (2026b). Diseño propuesto: LafroX (2026e, diseño de este apartado, pp.~\pageref*{fuente-19-anexo-4-1-p-tecnologias-soporte}--\pageref*{fuente-19-anexo-4-1-p-tecnologias-soporte-fin}).\par}

La aplicación Android de terreno soporta la versión mayor vigente de Android y las dos anteriores. Antes de cada liberación se verifican los flujos y periféricos en la matriz por modelo. Cada paquete Kotlin se firma y se distribuye con la gestión de flota de Android Enterprise y Zebra DNA. Antes de instalar o actualizar, se verifican la firma y la huella contra el artefacto aprobado y se rechazan los paquetes alterados o de origen no autorizado. Las actualizaciones siguen las tandas por sitio y enlace de la ventana dominical descrita en 4.2.6.6 del Subdocumento 4, sin alterar los plazos de corrección de vulnerabilidades ni sus restricciones de calendario (RT-17.03 y RT-17.04, Bases Técnicas Transversales, cap. 17, p. 31).

La Tabla~\ref{tab:inventario-software} incorpora la matriz de navegadores y versiones soportadas de RT-13.10 (Bases Técnicas Transversales, cap. 13, p. 26), aplicable a los portales y a la consola web de 4.1.3.1. La matriz utiliza versiones relativas para mantener el soporte durante el contrato. Antes de cada liberación se registra el número exacto, sistema operativo, resultado y fecha de las pruebas. Angular 22 usa el conjunto Baseline del 7 de mayo de 2026. Las versiones ofertadas deben pertenecer a ese conjunto y, al migrar Angular, al de la versión adoptada (Angular, 2026b). Cada nueva versión estable de navegador activa pruebas de autenticación, navegación, formularios, teclado y adaptación de pantallas en QA antes de sustituir la versión más antigua de la matriz. Un fallo genera corrección y repetición de pruebas antes de declarar su compatibilidad. Los parches de seguridad se comprueban y se incorporan a la matriz con prioridad. La revisión trimestral de soporte del framework se complementa con esta comprobación por cada versión estable del navegador.

El fin de soporte no publicado se registra como \emph{no publicado}, con responsable y frecuencia de revisión. No significa soporte indefinido. Antes de aprobar una dependencia se verifica licencia, mantenimiento, vulnerabilidades y alternativa de sustitución (RT-11.26). La ficha de liberación identifica versión exacta, EOL conocido o política aplicable, prueba, rollback y responsable. La aceptación rechaza una imagen fuera de soporte.

Las alternativas se evalúan por operación offline, integración industrial, transacciones, geografía, carga del equipo y reversibilidad. PostgreSQL/PostGIS se selecciona frente a MariaDB por el dominio espacial. Angular se selecciona por su organización en componentes y sus contratos TypeScript. Kotlin se selecciona frente a una capa móvil adicional por validación del parque. Keycloak mantiene identidad propia frente a delegar el dominio de autorizaciones a un IdP propietario. Las herramientas de construcción no cambian esos contratos. Los criterios concretos se verifican en ADR-04/06/07/20/21 y el Anexo 4-T.
\subsection*{Implementación del backend}

La tabla relaciona cada capacidad del backend con su implementación en Laravel/PHP y con la comprobación que verifica sus contratos.
\begin{tablalafrox}{Implementación lógica del backend Laravel}{tab:mapeo-laravel}{F{0.27} F{0.35} Y}{\cab{Capacidad} & \cab{Implementación Laravel/PHP} & \cab{Comprobación}}
APIs y módulos & Rutas y controladores Laravel, servicios por contexto PSR-4. & Contratos OpenAPI y límites de módulo. \\
Persistencia y geografía & PDO/Query Builder y SQL PostGIS parametrizado, migraciones Laravel basales. & Datos, índices y consultas espaciales verificados. \\
Tareas y planificación & Consumidor PHP de JSON en SQS FIFO, trabajos Laravel en colas separadas y planificador único. & Orden por grupo, reintentos y tareas únicas. \\
Cola local y shipper & Adaptador AMQP con \texttt{php-amqplib} y sobre JSON. & Corte de 24 h, confirmación y drenaje sin pérdida. \\
Identidad y administración & Guard OIDC, políticas por recurso y portal Angular. & Permisos, baja y relevos de turno. \\
Contratos, pruebas y trazas & \texttt{swagger-php}, AsyncAPI, PHPUnit y OpenTelemetry PHP. & Contratos de API y eventos verificados, con correlación. \\
\end{tablalafrox}

{\fontsize{9pt}{11pt}\selectfont\raggedright Fuente: elaboración a partir de Distribuidora Puelche S.A. (2026c, secciones~14.1--14.2 y~16.1, pp.~24--25 y~28--29), para volumetría y decisiones abiertas, y LafroX (2026e, diseño de este apartado, pp.~\pageref*{fuente-19-anexo-4-1-p-tecnologias-soporte}--\pageref*{fuente-19-anexo-4-1-p-tecnologias-soporte-fin}).\par}

\label{fuente-19-anexo-4-1-p-tecnologias-soporte-fin}
